% Revised abstract for ELAPSE paper (Physica A submission)
% Word count target: <=200 words

We establish that entropy-triggered data deletion on complex networks undergoes
a percolation-like phase transition at a critical threshold $H_c^*$, controlled
by the algebraic connectivity $\lambda_2$ of the network Laplacian.
The time-integrated entropy exposure (IEE), measuring cumulative privacy risk,
serves as the order parameter: IEE diverges algebraically as $H_c \to H_c^*$
from below, and scales as $\mathrm{IEE} \sim n^{\alpha}$ with $\alpha \approx
1.55$ across Erd\H{o}s--R\'enyi, Barab\'asi--Albert, and Watts--Strogatz
topologies.
We compare five deletion mechanisms---EGDM entropy-gate (M1), SIR epidemic
spread (M2), stochastic first-passage (M3), Hill-function cooperative mortality
(M4), and social cascade pressure (M5)---and an ensemble (M6) whose voting
weights are learned via entropy-regularised Nelder--Mead optimisation.
Theoretical guarantees cover the full stochastic system: well-posedness (Theorem~1),
a second-moment Gronwall IEE bound (Theorem~2), and spectral convergence of
the normalised distribution with It\^{o} correction (Corollary~1).
Real-world validation on SNAP Gnutella P2P and academic collaboration networks
confirms $\geq 9\times$ IEE reduction versus the no-deletion baseline.
M6 outperforms any single mechanism under topology uncertainty, on evolving
networks, and when the critical threshold $H_c$ is misspecified, establishing
the ensemble's value beyond the single-topology setting.
